% ============================================================
%  Contenu du CV — version française
%  Les macros \cvexp, \cvedu, \cvskill... sont définies par chaque
%  modèle (une colonne / deux colonnes) : on écrit le contenu une seule fois.
%
%  \cvexp{Titre}{Organisation}{Lieu}{Date}{ \cvli{puce} \cvli{puce} }
%  \cvedu{Titre}{Organisation}{Lieu}{Date}
%  Les variantes \cvexplast / \cvedulast s'utilisent pour le DERNIER
%  élément d'un bloc (elles terminent la frise du modèle deux colonnes).
% ============================================================

% ---------- Libellés ----------
\newcommand{\cvLabProfile}{Profil}
\newcommand{\cvLabExp}{Expériences}
\newcommand{\cvLabProj}{Projets \& open source}
\newcommand{\cvLabEdu}{Formation}
\newcommand{\cvLabEduAwards}{Formation \& distinctions}
\newcommand{\cvLabSkills}{Compétences}
\newcommand{\cvLabMisc}{Divers}
\newcommand{\cvLabNetworks}{Réseaux}
\newcommand{\cvLabLang}{Langues}
\newcommand{\cvLabAwards}{Distinctions}
\newcommand{\cvLabInterests}{Centres d'intérêt}
\newcommand{\cvColon}{\ :\ }

% ---------- Titre et profil ----------
\newcommand{\cvHeadline}{Élève-ingénieur EPITA \,·\, IA agentique \,·\, C++ \,·\, Cryptanalyse}
\newcommand{\cvTagline}{Informatique \& Recherche}
\newcommand{\cvProfile}{Élève-ingénieur à l'EPITA (diplôme prévu en 2028), stagiaire IA chez Opensee sur une solution d'IA agentique pour l'analyse de données financières et étudiant-chercheur en cryptanalyse différentielle (C++ parallèle). Contributeur open source (deux PR fusionnées dans \texttt{langchain-aws}). Je cherche à travailler là où l'ingénierie logicielle, l'IA et la recherche se rencontrent.}

% ---------- Expériences ----------
\newcommand{\cvExperiences}{%
  \cvexp{Stagiaire IA}{Opensee}{Puteaux}{Sept. 2026 -- présent}{%
    \cvli{Contribution au développement d'\textbf{Agensee}, solution d'IA agentique dédiée à l'analyse de données financières, au sein de l'équipe IA.}}
  \cvexp{Étudiant-chercheur, équipe Sécurité \& Systèmes}{Laboratoire de Recherche de l'EPITA (LRE)}{Lyon}{Depuis fév. 2026}{%
    \cvli{Optimisation de la minimisation des \textbf{S-boxes actives} dans l'outil de cryptanalyse différentielle \textbf{TAGADA} : algorithmes de programmation dynamique \textbf{multithreadés en C++}, adaptés aux contraintes de la bibliothèque.}
    \cvli{Automatisation de la récupération des tours de chiffrement (parsing JSON), benchmarks conteneurisés (Docker), veille scientifique et synthèse des avancées à l'équipe.}}
  \cvexp{Assistant pédagogique en programmation}{EPITA (cycle préparatoire) et SupBiotech}{Lyon}{2025 -- 2026}{%
    \cvli{Accompagnement quotidien d'étudiants de première année en \textbf{C} et \textbf{Python}, Git et Linux (NixOS) ; relecture des sujets et tests de TP, correction (Criterion, GitLab).}}
  \cvexplast{Bénévole}{Ensemble scolaire Charles de Foucauld}{Lyon}{2016 -- 2023}{%
    \cvli{Animation d'activités (escape game), tenue de caisse, référent de stand avec une équipe d'une vingtaine de personnes.}}}

% ---------- Projets ----------
\newcommand{\cvProjects}{%
  \cvexp{Contributions à \texttt{langchain-aws} (LangChain)}{2 PR fusionnées}{\href{https://github.com/langchain-ai/langchain-aws/pull/1389}{\#1389}, \href{https://github.com/langchain-ai/langchain-aws/pull/1415}{\#1415}}{Oct. 2026}{%
    \cvli{Ajout des niveaux d'effort de raisonnement (\texttt{reasoning\_effort\_levels}) aux profils de modèles Amazon Bedrock ; valeurs vérifiées à la source, lint et tests validés.}}
  \cvexp{Emulators Collection}{Projet personnel · C++20, WebAssembly, React}{\href{https://github.com/Vorpalin/emulators-collection}{GitHub}}{}{%
    \cvli{Émulateurs \textbf{CHIP-8, Super-CHIP, Atari 2600 et Game Boy} : cœurs C++20 compilés en \textbf{WebAssembly} (Emscripten), frontend React/TypeScript ; CI GitHub Actions, analyse de sécurité, releases automatiques.}}
  \cvexp{PING, IDE en ligne pour débutants}{EPITA · Java, React}{Équipe agile de 5}{Depuis mars 2026}{%
    \cvli{Conception produit (état de l'art, interviews, maquettes), interface React, backend Java testé par pipeline CI/CD GitLab, documentation Javadoc.}}
  \cvexplast{42sh, shell POSIX en C}{EPITA}{Équipe agile de 4}{Janv. 2026}{%
    \cvli{Lexer et parser, architecture fondée sur les patrons \textbf{visiteur} et \textbf{composite}, revue des merge requests GitLab. Autre projet : Yakamon, API REST Java en couches avec JUnit.}}}

% ---------- Formation ----------
\newcommand{\cvEducation}{%
  \cvedu{Diplôme d'ingénieur en informatique}{EPITA}{Lyon}{Sept. 2023 -- 2028}
  \cvedu{Semestre à l'étranger, Computer Science}{Centria University of Applied Sciences}{Kokkola, Finlande}{Jan. -- mai 2025}
  \cvedulast{Baccalauréat général}{Lycée Charles de Foucauld}{Lyon}{Juil. 2023}}

\newcommand{\cvAwards}{%
  \cvaward{Top Scorer 2026}{Challenge du Meilleur étudiant en informatique de France}
  \cvaward{Open source}{2 PR fusionnées dans langchain-aws}}

% ---------- Compétences, langues, intérêts ----------
\newcommand{\cvSkills}{%
  \cvskill{Langages}{C, C++ (C++20), Python, Java, C\#, JavaScript / TypeScript, SQL}
  \cvskill{IA \& recherche}{IA agentique, cryptanalyse différentielle, programmation dynamique, calcul parallèle}
  \cvskill{Web \& WASM}{React, Vite, WebAssembly (Emscripten), HTML/CSS, API REST}
  \cvskill{Outils}{Git (GitHub, GitLab), CI/CD, Docker, CMake, Linux (NixOS), JUnit, Criterion}
  \cvskill{Bases de données}{PostgreSQL, MySQL}}

\newcommand{\cvLangItems}{%
  \cvlangitem{Français}{Langue maternelle}
  \cvlangitem{Anglais}{Niveau C1 · TOEIC 960/990}
  \cvlangitem{Italien}{Notions}}
\newcommand{\cvLangShort}{Français (natif), Anglais (C1, TOEIC 960/990), Italien (notions)}

\newcommand{\cvInterests}{%
  \cvinterest{Tennis de table}{Départemental 1 en équipe, pré-régional en individuel}
  \cvinterest{Cinéma}{}}
